\documentclass[a4paper,11pt]{article}
% Package to make citations superscrit with brackets
\usepackage[super,square]{natbib}
% Package to change margin size
\usepackage{anysize}
\marginsize{1.5cm}{1.5cm}{1cm}{1.5cm} %left/right/top/bottom
% Package to make headers
\usepackage{fancyhdr}
\renewcommand{\headrulewidth}{0pt}
% Package for highlights
\usepackage{soul}
% Colors for the references links
\usepackage[dvipsnames]{xcolor}
% Package to link references
\usepackage{hyperref}
\usepackage[]{graphicx}
\usepackage{float}
\usepackage{physics}
\usepackage{tensor}
\usepackage{esint}
\usepackage{amsfonts}
\usepackage{mathtools}
\usepackage{amsmath}
\usepackage{amssymb}
\usepackage{enumitem}
\usepackage{setspace} %to set line spacing, see TOC for an example
\usepackage[english]{babel}
\hypersetup{
    colorlinks=true,
    linkcolor=black,
    citecolor=CadetBlue,
    filecolor=CadetBlue,      
    urlcolor=CadetBlue,
}
% Package for lorem ipsum
\usepackage{lipsum}
% Package for multicolumn
\usepackage{multicol}

\usepackage{multicol,caption}
% Package for removing paragraph identations
\usepackage{parskip}
\setlength\columnsep{18pt} 
\setlocalecaption{english}{figure}{Fig.}
\setlocalecaption{english}{table}{Tab.}
\counterwithin{figure}{section} % comandi per numerazione interna di figure, tabelle ed equazioni
\counterwithin{table}{section}
\counterwithin{equation}{section}
\setcounter{tocdepth}{2}


%Defining Chapter, a section bigger than \section
\newcounter{Chapcounter}
\newcommand\showmycounter{\addtocounter{Chapcounter}{1}\themycounter}
\newcommand{\chapter}[1] 
{ {\centering          
    \addtocounter{Chapcounter}{1} \huge \underline{\textbf{#1}} } 
    %\addcontentsline{toc}{section}{\underline{#1}}    
    \addcontentsline{toc}{section}{\large \underline{#1}}    
}
\newcommand\pd{\partial}
\newcommand{\COMMENT}[1]{} %comment block
 
\begin{document}
\begin{center}
\Large{\textbf{Introduction to Quantum Field Theory,\\ Lecture notes}}
\vspace{0.4cm}
\normalsize
\\ Università degli Studi di Torino\\Academic Year 2026/2027 \\Prof. Boglione,
notes by Paolo Rocchietti March\footnote{paolo.rocchiettimarc@edu.unito.it} \\

\vspace{0.1cm}
%\textit{Organization}
\medskip
\normalsize
\end{center}
\vspace{0.4cm}
{\color{gray}\hrule}
\null\vfill
%\begin{figure}[H]
%    \centering
%\includegraphics[width = 0.8\linewidth]{./frontpage.png}
%\end{figure}
\vfill\null
%non mi convince, magari la cambierò
\newpage
\begin{spacing}{.9}
\tableofcontents
\end{spacing}
\newpage
\pagestyle{fancy}
\thispagestyle{fancy}
\fancyhead[L]{\small{\textbf{Università degli studi di Torino} \\ a.y. 2026/2027}}
\fancyhead[R]{\small{\textbf{Introduction to QFT notes}\\ Prof. Boglione}}
\fancyfoot[C]{\small{\thepage}}
%\input{./somefile.tex}
\section{Lesson 1, 23/09}
\textbf{PRIMA LEZIONE SU TABLET}

%-----------------------------------------------------------------------------------------------------
\section{Lesson 1, 23/09}
\textbf{PRIMA LEZIONE SU TABLET}
\section{Lesson 2, 24/09}
\subsection{Natural units}
From this point onwards, natural units will be used:
$$\hbar = c = 1$$
\textbf{Note:} this messes up dimensionality, but, convenientely, forces mass and energy to have
the same units
$$E=mc^2\underset{c=1}\rightarrow E=m,$$
and lenght and time to have the same dimensionality
$$x=ct\underset{c=1}\rightarrow[x]=[t].$$
To convert back into more familiar units the constants are, respectively,
$$c=299.792.458\text{ m/s}$$
$$\hbar c \approx 0.97 \text{ GeV$\cdot$fm}$$
\subsection{Klein-Gordon Equation}
The relativistic quantum mechanics' fundamental equation is
$$(\Box+m^2)\psi(x)=0$$
where
$$\Box = \partial^\mu\partial_\mu=\pdv[2]{}{t}-\laplacian$$
and $\psi(x)$ is still the Schrödinger's wavefunction.

It will now be shown that this equation \textbf{does not} conserve probability, using a continuity
equation.

Regular quantum mechanics states
$$\pdv[]{\rho}{t}+\div\mathbf{j} = 0,$$
where
$$\rho=\bra{\psi_t}\ket{\psi_t}\quad\quad \mathbf{j=}\frac{i}{m}\left(\psi^*\grad\psi-\psi\grad\psi^*
\right)$$

Consider the Klein-Gordon equation left-multiplied by $\psi^*(x)$ and its conjugate
$$\begin{dcases}
  \psi^*(x)\left(\Box+m^2\right)\psi(x)=0\\
  \psi(x)\left(\Box+m^2\right)\psi^*(x)=0
  \end{dcases}$$
their difference is
$$\psi^*(\Box+m^2)\psi - \psi(\Box+m^2)\psi^*=0$$
$$\partial^\mu(\psi^*\partial_\mu\psi - \psi\partial_\mu\psi^*)=0$$
$$\partial^\mu J_\mu=0$$
this is a 4D continuity equation and $J^\mu$ is a Noether current.

Expanding the equation yields
$$\pdv[]{}{t}\left(\psi^*\pdv[]{\psi}{t}-\psi\pdv[]{\psi^*}{t}\right)-\div\left(\psi^*\grad \psi
-\psi\grad\psi^*\right)=0$$
the second term is proportional to the quantum probability current, but it is missing an $i$!
Moreover, the first term, representing $\rho$, is \textbf{not definite positive}, so it cannot
represent the norm of the state $\ket{\psi_t}$!\\

This proves that it is \textbf{impossible} to make QM relativistic using wavefunctions and retaining
the probabilistic approach.

Paul Dirac used another approach, he wanted to decrease the order of space derivative (instead of
increasing the time derivative order, like K-G): this result was better, but for free particles
it yielded a double solution with $E<0$ states, violating a fundamental QM theorem!\\
For this reason, he devised the concept of the \textit{Dirac sea}:
these $E<0$ states must all be occupied (because they cannot be observed), and a uniform charge
can be factored out of Maxwell's equations with proper boundary conditions; when one of these
particles in the sea is excited, the "hole" left behind can be percieved as a "positive electron":
he predicted the existence of \textit{antiparticles}.

This interpretation is still classical, but it begins to convey the concept that a finite number
of particles is \textit{not} possible for relativistic QM, so the concept of wavefunctions must be
left behind in favor of field theories.
\subsection{Symmetries and conservations}
To build a theory, it is necessary to start with principles and symmetry:
\textbf{Lorentz-invariance} to guarantee the conservation of distance between points (this is also
useful to define a metric topology!).

Why symmetry? Noether theorem states that for every 4-current that respects a continuity equation
$$\partial_\mu J^\mu=0$$
there is a conserved Noether charge
$$Q:=\int_{}^{}\dd^3xJ^0\text{, so that }\dv[]{Q}{t}=0$$
\subsubsection{Lorentz-Poincaré group}
The Lorentz invariant is
$$s^2=t^2-x_ix^i = t'^2-x_i'x'^i = x_\mu x^\mu:$$
the Lorentz group is the set of all transformations $x^\mu\rightarrow x'^\mu$ over the Minkowski
space that conserve $s^2$; in matricial representation the elements are
$$x'^\mu = \Lambda\indices{^\mu_\nu} x^\nu$$
the invariant is conserved \textbf{if and only if}
$$\boxed{g_{\mu\nu}\Lambda\indices{^\mu_\rho}\Lambda\indices{^\nu_\sigma} = g_{\rho\sigma}.}$$
The covariant quantities transform with the inverse matrix
$$x'_\mu = \Lambda\indices{_\nu^\mu} x'_\nu$$
\textbf{Observation 1} Since $g_{\mu\nu}$ is a symmetric matrix, this equation encompasses $10$
constraints on $\Lambda\indices{^\mu_\nu}$, leaving $3$ degrees of freedom for rotations and $3$ for
boosts.

\textbf{Naming convention} Consider the determinant of the invariance condition:
$$\det(\Lambda^Tg\Lambda)=\det(g)\leftrightarrow\det(\Lambda)^2=1\leftrightarrow\det(\Lambda)=\pm1$$
Transformations with $\det(\Lambda) = 1$ are called \textit{proper Lorentz transformations}.\\
Transformations with $\det(\Lambda)=-1$ are called \textit{imporper Lorentz transformations}.

\textbf{Naming convention} Consider the (0,0) component of the invariance condition:
$$g_{\mu\nu}{\Lambda^\mu}_0{\Lambda^\nu}_0 = 1 \leftrightarrow (\Lambda_{00})^2 - \Lambda\indices{^i_0}
g_{ii}{\Lambda^i}_0=1\leftrightarrow |{\Lambda^0}_{0}|^2\geq1$$
Because this term regulates how the transformation relates $x^0$ and $x'^0$,
transformations with $\Lambda\indices{^0_0}\geq1$ are called
\textit{orthocronous Lorentz transformations}, while those with $\Lambda\indices{^0_0}\leq1$ are
called \textit{non-orthcoronous Lorentz transformations}.

These 2 divisions split the Lorentz group into 4 parts
\begin{itemize}
    \item Proper and orthocronous $L^{\uparrow}_{+}$: restricted Lorentz group, sometimes referred to
	as "Lorentz Group";
    \item Improper and orthocronous $L^{\uparrow}_-$: spacial inversion;
    \item Proper and non-orthocronous $L^{\downarrow}_+$ : time inversion;
    \item Improper and non-orthocronous $L^{\downarrow}_-$: total inversion of spacetime.
\end{itemize}
\textbf{Note} The only subgroup is $L^\uparrow_+$, as it is the only one that contains the identity
matrix $\mathbb{I}$.

%-------------------------------------------------------------------------------------
\section{Lesson 3, 28/09}

\section{Lesson 4, 30/09}
The generators are useful because one can study the entire group just by considering
the generators: they are operators uniquely defined by their commutation relations
using Lie brackets; the set of generators equipped with their commutation relations
is the \textbf{Lie algebra} associated with the group.

Once the algebra is known, the choice of the representation yields matrices equivalent
to the generators.

As first examples, it is useful to consider $SO(3)$ and $SU(2)$, because they are
closely related to the Lorentz group (which is the final purpose of this parenthesis
on group theory).
\subsection{Example: angular momentum as a representation}
As is known, the algebra of the angular momentum operators $J^i$ is
\begin{gather*}
    \left[J^i,J^j\right] = i\varepsilon\indices{^i^j_k}J^k
\end{gather*}
the rank of a group is the maximum number of generators that commute with one another:
in this case the rank is $1$ and the chosen operator is $J^3$.\\
A Casimir operator is an operator that commutes with all generators: in this case it
is
$$J^2 = J_iJ^i.$$
The physical question is then "how many possible values/eigenstate of the angular
momentum are possible?" and it turns out that the answer to this question also
answers, in a QFT formalism, "how many particles does the field predict?".\\
The process to find the possible eigenvalues and eigenstates is known (see QM1 course)
and it yields the fact that for each $j$ value, the number of eigenstates of $J^3$
is finite and their eigenvalues $m$ must be in the range
$$-j< m < j$$
where $j$ is called the \textit{maximum weight}; so for each $j$ the possible values
of $J^3$ (and so the number of different eigenstates) are $2j+1$. From this process
also follows that the only possible values of $j$ are positive integers or
semi-integers and that the values of $m$ always differ by an integer (so if $j$ is
semi-integer, all the corresponding values of $m$ must also be semi-integer and
vice-versa).

This formalism is very powerful because it relates the irreducible representations
of the algebra to its vector basis: in a QFT, the dimension of the vector space
will correspond to the number of possible physical particles.

This algebra is common between $SU(2)$ and $SO(3)$, with the key difference that
for $SO(3)$ only integer values of $j$ are permitted: as seen in intro groups, $SU(2)$
can be represented as a 3-sphere immersed in $\mathbb{R}^4$, while $SO(3)$ is
only half 3-sphere in which antipodal points are mapped to the equator (??).

\subsection{Example: reducible representation}
The combination of 2 spin $\frac{1}{2}$ particles yields a reducible representation:
the vector space is
$$V = \mathbb{C}^2\otimes\mathbb{C}^2,\quad \dim(V)=4$$
again, the process of spin combination is treated in the QM1 and QM2 courses and it
yields a \textit{triplet} and a \textit{singlet} of states:
\begin{align*}
&\text{Triplet (J=1): }\begin{dcases*}
\ket{J=1,M=-1}\\
\ket{J=1, M=0}\\
\ket{J=1, M=-1}
\end{dcases*}\rightarrow\text{Vector states}\\
\phantom{x}\\
&\text{Singlet J=0}: \ket{J=0,M=0} \rightarrow\text{Scalar state}
\end{align*}
So the space $V$ is subdivided into two different vector spaces which differ in
symmetry: the singlet is antisymmetric while the triplet is symmetric; this
representation is reducible, because the image of the representation of the algebra
is the direct sum of the images of the representations of two different algebras.
\subsection{The Lorentz Group Representation}
Only the subgroup $L^\uparrow_+$ is here considered and its formal name is
$SO(1,3)$, where $(1,3)$ is the signature of the Minkowski metric $(+,-,-,-)$.

Consider an infinitesimal Lorentz transformation
$$\Lambda\indices{^\mu_\nu} = \delta\indices{^\mu_\nu}+\omega\indices{^\mu_\nu},\quad
|\omega\indices{^\mu_\nu}|<<1\,\forall\mu,\nu
$$
where $\omega\indices{^\mu_\nu}$ are called \textit{parameters} of the transformation.

Substituting into the invariance condition yields
\begin{gather*}
g_{\mu\nu} \Lambda\indices{^\mu_\rho}\Lambda\indices{^\nu_\sigma} = g_{\rho\sigma}\\
g_{\mu\nu} \left(\delta\indices{^\mu_\rho}+\omega\indices{^\mu_\rho}\right)
\left(\delta\indices{^\nu_\sigma}+\omega\indices{^\nu_\sigma}\right)=g_{\rho\sigma}\\
g_{\mu\nu}\delta\indices{^\mu_\rho}\delta\indices{^\nu_\sigma} + 
\omega_{\nu\rho}\delta\indices{^\nu_\sigma} +
\omega_{\mu\sigma}\delta\indices{^\mu_\rho}=g_{\rho\sigma}\\
\omega_{\rho\sigma} = -\omega_{\rho\sigma}
\end{gather*}
so the parameters must be \textbf{antisymmetric}. This confirms the fact that
each transformation is uniquely determined by 6 quantities.
The transformation becomes
$$\Lambda\indices{^\mu_\nu}  =\delta\indices{^\mu_\nu} +
g^{\mu\rho}g\indices{^\sigma_\nu}\omega_{\rho\sigma}$$
since $\omega_{\rho\sigma}$ are antisymmetric, the only non-zero part of the product
$\omega_{\rho\sigma}g^{\mu\rho}g\indices{^\sigma_\nu}$ is given by the antisymmetric
part of $g^{\mu\rho}g\indices{^\sigma_\nu}$: anti-symmetrizing it yields
$$\Lambda\indices{^\mu_\nu} = \delta\indices{^\mu_\nu}+\frac{1}{2}\omega_{\rho\sigma}
\left(g^{\mu\rho}g\indices{^\sigma_\nu}-g^{\mu\sigma}g\indices{^\rho_\nu}\right)=
\delta\indices{^\mu_\nu} + i\frac{1}{2}\omega_{\rho\sigma}\left(M^{\rho\sigma}\right)
\indices{^\mu_\nu}
$$
confronting with the definition of generators
$$g=\mathbb{I} + i\alpha_aT^a$$
yields that the generators are
$$M^{\rho\sigma} = -i\left(g^{\mu\rho}g\indices{^\sigma_\nu}-g^{\mu\sigma}
g\indices{^\rho_\nu}\right)$$
so the generic Lorentz transformation can be expressed as
$$\Lambda = \exp\left[\frac{i}{2}\omega_{\rho\sigma}M^{\rho\sigma}\right]$$
and its commutation relations are
$$\left[M^{\lambda\tau},M^{\rho\sigma}\right]= -i \left(
    g^{\lambda\sigma}M^{\tau\rho}+g^{\tau\rho}M^{\lambda\sigma}-g^{\lambda\rho}
    M^{\tau\sigma}-g^{\tau\sigma}M^{\lambda\rho}
\right)$$
To make computations easier, the 6 generators are split into $3+3$ easier operators:
since $SO(3)$ is a subgroup of $SO(1,3)$, there is a way to express half of the
Lorentz Group generators in terms of those of $SO(3)$:
$$J_i  =\frac{1}{2}\varepsilon_{ijk}M^{jk}$$
which obey the $SO(3)$ algebra
$$\left[J_i,J_j\right] = i \varepsilon\indices{_i_j^k}J_k$$
and generate rotations; the boosts are generated by
$$K^i = M^{0i}$$
which obey the commutation relation
$$\left[K^i,K^j\right] = -i \varepsilon\indices{^i^j_k}J^k$$
which means that the algebra of the $K$'s is \textbf{not} closed, whereas the algebra
of the $J$'s is: this is due to the fact that 3D rotations form a subgroup of
$SO(1,3)$, while boosts do not (it is easy to check that the composition of 2 Cartesian
boosts is the composition of a Cartesian boost and a rotation).
The mixed commutator yields
$$[J_i,K_j] = i\varepsilon_{ijk}J^k$$
so the set of all $J$'s and all $K$'s is the set of the generators of $SO(1,3)$.

\newpage
\section{Lesson 5, 01/10}
\subsection{More on the Lorentz group}
The generic Lorentz transformation can also be written in terms of $J,K$:
$$\Lambda = \exp\left[i\vartheta_iJ^i+i\eta_iK^i\right]$$
where $\vartheta_i$ are the Euler angles (parameters of the rotation) and $\eta_i$
are the rapidities (parameters of the boost).

With some work, it is possible to prove that
$$\frac{1}{2}M_{\mu\nu}M^{\mu\nu} = J_iJ^i - K_iK^i = \mathbf{J^2}-\mathbf{K^2}$$
and
$$-\frac{1}{4} \varepsilon^{\mu\nu\rho\sigma}M_{\mu\nu}M_{\rho\sigma} = J_iK^i = 
\mathbf{J}\cdot\mathbf{K}$$
are Casimir operators for the representation: there exist only 2 linear independent
Casimir operators and each and any linear combination of these two is valid, but
these two are important because they are, respectively, symmetric and antisymmetric
(how can they be?? They're scalar!).

\textbf{Theorem:} The algebra of the Lorentz Group can be factorized as the direct
product of two $SU(2)$ algebras
$$\mathfrak{so}(1,3) = \mathfrak{su}(2)\otimes\mathfrak{su}(2)$$

This suggests to find new generators that are suited to this decomposition
\begin{gather*}
\mathbf{N} = \frac{1}{2}\left(\mathbf{J}+i\mathbf{K}\right)\\
\mathbf{N}^\dagger = \frac{1}{2}\left(\mathbf{J}-i\mathbf{K}\right)
\end{gather*}
their algebra is
\begin{align*}
&\left[N_i,N_j^\dagger\right] = 0\\
&\left[N_i,N_j\right] = i\varepsilon_{ijk}N^k\\
&\left[N_i^\dagger,N_j^\dagger\right] = i\varepsilon_{ijk}N^{k\dagger}
\end{align*}
\textbf{Observation:} These generators are convenient, since the algebra
of the $N$'s and that of the $N^\dagger$'s are closed and disjoint.

The Casimir operators are
\begin{align*}
    N_iN^i &\longrightarrow n(n+1)\\
	   &\text{\footnotesize eigenvalues}\\
    N^\dagger_iN^{i\dagger} &\longrightarrow m(m+1)    ,
\end{align*}
so the irreducible representations of the group will be labeled by the pair $(n,m)$.
Since they obey $SU(2)$ algebra, it is clear that $n,m$ will have integer or
semi-integer values and that the eigenvalues of
$N_3,N_3^{\dagger}$, $s$ and $s^\dagger$, will respect
$$-n<s<n\qquad-m<s^\dagger<m.$$
To each pair $n,m$ corresponds a name for the field
\begin{table}[ht!]
    \centering
    \begin{tabular}{c|c|c|c}
	$(n,m)$&Field name&Dim&Example\\ \hline
	$(0,0)$&Scalar $\varphi$&1&Higgs \\
	$(1/2,0)$&Spinor $\psi_L$&2&Weyl Spinor\\
	$(0,1/2)$&Spinor $\psi_R$&2&Weyl Spinor\\
	$(1/2,0)\oplus(0,1/2)$&Dirac Spinor&4&\\
	$(1/2,1/2)$&4-vector field (spin 1)&4&$A^\mu$ ELM\\
	$(1,0)\oplus(0,1)$&Rank 2 antisymm tensor&6&$F^{\mu\nu}$ ELM
    \end{tabular}
\end{table}
\newpage
\subsection{Poincaré Group}
The QFT must also be invariant over spacetime \textit{translations}: the Poincaré
group includes Lorentz transformations \textit{and} translations
$$x^\mu\rightarrow x'^\mu:=\Lambda\indices{^\mu_\nu}x^\nu + a^\mu$$
the group parameters are $6+4=10$. The generators of the 4-translations,
similarly to QM, are the 4-momenta
$$P^\mu = (H,P^i);$$
moreover, in QM1 the time evolution operator was found to be
$$U=e^{-iH}$$
which, in hindsight, is an exponential map, in which $H$ is the generator and $t$ is
the parameter!

There will be some problems with the fact that the translation group is not
compact.

Finding the commutation relations for the algebra is not trivial, since rotations
and translations do not commute, even in Minkowski space. Consider two
successive Poincaré transformations (why????):
$$x^\mu\rightarrow{(\Lambda_{(1)})}\indices{^\mu_\nu}x^\nu + a_1^\mu \rightarrow
{(\Lambda_{(2)})}\indices{^\mu_\rho}{(\Lambda_{(1)})}\indices{^\mu_\nu}x^\nu + 
{(\Lambda_{(2)})}\indices{^\mu_\rho}a_1^\rho+a_2^\mu$$
The commutation relations are
\begin{align*}
&\left[P^\mu,P^\nu\right]=0\\
&\left[M_{\mu\nu},P_\rho\right] = -ig_{\rho\sigma}P_\nu+ig_{\nu\rho}P_\mu\\
\end{align*}
They are clearer when time and space indexes are expressed separately
\begin{align*}
&\left[P^i,P^j\right]=0\\
&\left[P^i,H\right]=0\\
&\left[K^i,P^j\right]=i\delta^{ij}H\\
&\left[J^i,J^j\right]=i\varepsilon\indices{^i^j_k}J^k\\
&\left[K^i,H\right]=iP^i\\
&\left[J^i,P^j\right]=i\varepsilon\indices{^i^j_k}P^k\\
&\left[J^i,H\right]=0\\
&\left[K^i,K^j\right]=i\varepsilon\indices{^i^j_k}J^k\\
&\left[J^i,K^i\right]=i\varepsilon\indices{^i^j_k}K^k\\
\text{REORDER}
\end{align*}
Next, the Casimir operators: a natural choice is
$$P^2 = P_\mu P^\mu$$
a not so natural choice is the \textit{Pauli-Lubanski vector} squared:
$$W^\mu = \varepsilon^{\mu\nu\rho\sigma}P_\nu M_{\rho\sigma}\rightarrow W^2 = W_\mu
W^\mu.$$
A nice relation is
$$P_\mu W^\mu = 0$$
for antisymmetry reasons.


\section{Lesson 6, 05/10}
\subsection{Massive and massless particles representations}
\textbf{Massive particles}\\
It is known from SR that if a particle is massive, there exist a reference frame
in which the particle is at rest, i.e.
$$P^\mu = (m,\mathbf{0})$$
which means that
$$W^0 = 0\quad W^i = -\frac{1}{2} \varepsilon^{0ijk}M_{jk}$$
recalling the definition of $J^i$ this is
$$W^i = -\frac{1}{2}mJ^i$$
since the particle is at rest, this cannot represent an orbital angular momentum,
so it must be intrinsic, i.e. \textit{spin}:
$$W^i = -\frac{1}{2}mS^i$$
and so massive particles are identified by Casimir operators $P^2$ and $S^2\propto
W^2$ and by the commuting operators $S_3$ and $P^\mu$:
there are, for each value of $m,s$, $2s+1$ internal degrees of freedom.

\textbf{Massless particles}\\
If the particle is massless,
$$P_\mu P^\mu = m^2 = 0$$
which means
$$W_\mu W^\mu = 0,$$
and in general
$$P_\mu W^\mu = 0:$$
from these 3 conditions it can be shown\footnote{moodle?} that
$$W^\mu = h P^\mu$$
where $h = \pm 1$ is called \textbf{helicity} and represents the projection of the
angular momentum on the direction of the momentum
$$h = \frac{\mathbf{J}\cdot\mathbf{P}}{|\mathbf{P}|};$$
in this case, they can only be parallel or anti-parallel: these states are called
\textit{polarizations} and, for photons, they correspond to light polarizations.

This framework also allows a sector of the theory in which spin is continuous
($m=0$, $W^2<0$ states are like this, but they have not been observed in nature)
and a sector in which some particles, tachyons can move at $v>c$.

The key takeaway from this arguments is that

\noindent\fbox{%
    \parbox{\textwidth}{%
        \textbf{The irreducible representations of the Poincaré Group are in
        one-to-one correspondence with the possible types of elementary particles.\\
        The mathematical structure of space-time symmetries determines what kinds
        of particles can exist in nature.}
    }%
}
\subsection{Non-finite representations of groups: field representations}
Until this point, all the representations have been finite-dimensional, with
generators represented as matrices, but a field is continuous.

\textbf{Def.} A field is a well-behaved function of spacetime
$$\phi_a(x^\mu) = \begin{pmatrix}
    \phi_1(x^\mu)\\
    \phi_2(x^\mu)\\
    ...
\end{pmatrix}
$$
matrices can transform a field into another with linear combinations of the old
components as new components: these transformations are finite-dimensional, while
the ones on the coordinates $x^\mu$ can also be infinite-dimensional.

In QM the allowed transformations were unitary transformations, as they preserve
squared norms, which are the observable part of the theory: these transformations
are infinite-dimensional (think of the time evolution: its generator is the
Hamiltonian $H$, which is a differential operator).

In parallel to the definition of the angular momentum operator in QM, the Lorentz
generator for infinite-dimensional transformation is
$$M_{\mu\nu} = i (x_\mu\partial_\nu-x_\nu\partial_\mu)$$
and so a field transforms like
$$\phi_a(x) \rightarrow \phi_a'(x') = M\indices{^a_b}(\Lambda)\phi_b(\Lambda x)$$
where $\Lambda$ are infinite-dimensional transformations and $M\indices{^a_b}$ are
finite-dimensional;
the total transformation can be expressed with the exponential map
$$\phi_a(x')=\exp\left[-\frac{i}{2}\omega^{\rho\sigma}M_{\rho\sigma}\right]\phi_b(x),
$$
with
$$M_{\mu\nu} = M_{\mu\nu}^\text{inf} + M_{\mu\nu}^\text{fin}:$$
where the finite part generates the transformations of the field components\footnote{
Q: why are we assuming that the parameters are the same number and are the same?}
and the infinite part those of the space-time coordinates. This generators can be
expressed together thanks to the BCH formula.

\section{Lesson 7, 06/10}
\subsection{Classical Field theory}
Before delving into QFT, it is necessary to understand classical field theory.

A classical system is described by the generalized coordinates
$$q_i(t),\;\text{with }i=1,2...N$$
where $N$ is the (finite) number of degrees of freedom and $t$ is a parameter.
From this the Lagrangian is constructed
$$L(q,\dot q).$$
A field theory replaces the set of the coordinates$\left\{q_i\right\}$ with a
\textit{continuous function of space-time}, called a \textbf{field}:
$$\phi(t,\mathbf{x}),\;\mathbf{x}\in\mathbb{R}^3$$
the continuous coordinates $\mathbf{x}$ replace the discrete index $i$ and so the
degrees of freedom become infinite.

The next step is to construct a Lagrangian density
$$\mathcal{L}(\phi,\pd_\mu\phi)$$
and solve for the equations of motion.

\textbf{Example:} from finite to infinite degrees of freedom.\\
Suppose a system of $N$ point-like masses interacting with an elastic potential:
the Lagrangian will be
$$L=\sum_{k=1}^{N}\left[\frac{1}{2}m_i\dot x_i^2-\frac{1}{2}k(x_i-x_{i+1})^2\right].$$
Consider $m_i = m$ and the distance between them to be constant and equal to
$\varepsilon$:
$$L = \sum_{k=1}^{N}\frac{\varepsilon}{2}\left[\frac{m}{\varepsilon}\dot x_i^2
    -k\frac{(x_i-x_{i+i})^2}{\varepsilon}\right]$$
The limit as $\varepsilon$ goes to 0 is
$$\lim\limits_{\varepsilon\to0} L = \int_{}^{}\dd x \frac{1}{2}\left[
\mu\dot\phi(x)^2 - Y\left(\pdv[]{\phi}{x}\right)^2\right]$$
where $\frac{m}{\varepsilon}\rightarrow\mu$, $k\varepsilon\rightarrow Y$
(which is a Young module) and $x_i(t)\rightarrow\phi(t,x)$.
Applying the least action principle yields the Euler-Lagrange equation
$$\pdv[2]{\phi}{x} - \frac{\mu}{Y}\pdv[2]{\phi}{t} = 0.$$\\

In general, the Lagrangian density is a function of the spacetime coordinates $x^\mu$
through the dependency on a field and its derivatives
$$\mathcal{L}(x^\mu) = \mathcal{L}(\phi,\pd_mu\phi)$$
and the Lagrangian is its integral on the physical space
$$L(t) = \int_{}^{}\dd^3 x \mathcal{L}(\phi,\pd_\mu\phi);$$
another important quantity is the action, i.e. the integral over spacetime of the
density:
$$S = \int_{}^{}\dd^4x \mathcal{L}(\phi,\pd_\mu\phi).$$\\

All the theories constructed must be \textit{local}, i.e. they must obey causality,
so in the Lagrangian density $\mathcal{L}$ will never present terms like
$$\phi(x^\mu)\phi(y^\mu),\;\text{with }x^\mu\neq y^\mu$$
which represents an instantaneous interaction between two points of spacetime.

\textbf{Def.} A field theory is local when the dependency con $x^\mu$ is
$$\mathcal{L}(x^\mu) =\mathcal{L}(\phi(x^\mu),\pd_\mu\phi(x^\mu))$$
which means that the dynamics of the system do not connect different points in space
instantaneously.

\subsubsection{Transformations of fields}
Since the theory must be covariant, it is pivotal to understand how fields
behave under Poincaré transformations: consider a point in spacetime expressed in
coordinates as $x^\mu$ and the value of a field at that point $f(x^\mu)$; consider
then an infinitesimal transformation $x'^\mu = x^\mu+\delta x^\mu$ so that
$$\Delta f = f'(x'^\mu) - f(x^\mu) = f'(x^\mu+\delta x^\mu) - f(x^\mu) = f'(x^\mu)
-f(x^\mu) + \pd_\mu f'(x^\mu)\delta x^\mu + O(\delta^2 x^\mu)
$$
at first order, noting that $\pd_\mu f'(x^\mu) = \pd_\mu f(x^\mu) + O(\delta x^\mu)$,
this can be written as
$$\Delta f = \delta f +\pd_\mu f(x^\mu)\delta x^\mu.$$
This can be seen as the definition of the operator $\Delta$ on fields
$$\Delta = \delta + \delta x^\mu \pd_\mu:$$
under a Poincaré transformation, the variation of a field is made up of the
functional variation in the field itself ($\delta$) and the variation in the field
given by the transformation of the coordinates (called transport term).

\textbf{Example:} spacetime translations.\\
Consider the infinitesimal Poincaré transformation
$$x^\mu\rightarrow x'^\mu = x^\mu+\varepsilon^\mu$$
by definition, a local field $f$ is invariant under translations, so
$$\Delta f = 0$$
and this means
$$\delta f = - \varepsilon^\mu\pd_\mu f$$
and since $P_\mu = i\pd_\mu$, 
$$\delta f =-i\varepsilon^\mu P_\mu f.$$
This means that this functional transformation of the field is generated by the
4-momenta
$$f' = \exp\left[-i\varepsilon^\mu P_\mu\right]f$$

The Lorentz transformations, on the other hand, depend on the representation the field
belongs to: the simplest is the scalar representation.

\textbf{Example:} Lorentz transformation of a scalar field.\\
A scalar field is a field invariant under Poincaré transformations
$$\phi(x^\mu) = \phi'(x'^\mu).$$
As seen before, the translations are easy
$$\phi(x^\mu)\rightarrow\phi'(x'^\mu) = \exp\left[-i\varepsilon^\mu P_\mu\right]
\phi(x^\mu).$$
Consider now a generic Lorentz transformation
$$x^\mu\rightarrow x'^\mu = x^\mu + \varepsilon^{\mu\rho}x_\rho$$
where $\varepsilon^{\mu\rho}$ are the 6 completely antisymmetric parameters of the
transformation.\\
By definition,
$$\Delta \phi = 0 = \delta\phi + \varepsilon^{\mu\rho}x_\rho\pd_\mu \phi$$
and so
$$\delta\phi = -\varepsilon^{\mu\rho}x_\rho\pd_\mu\phi \rightarrow \delta
=-\varepsilon^{\mu\rho}x_\rho\pd_\mu$$
this operator is non-zero only for the antisymmetric part of $x_\rho\pd_\mu$, so
antisymmetrizing yields
$$\delta = -\varepsilon^{\mu\rho}\frac{1}{2}\left(x_\rho\pd_\mu-x_\mu\pd_\rho\right)=
-\varepsilon^{\mu\rho}\frac{i}{2}L_{\rho\mu}
$$
where $L_{\rho\mu}$ is the generalized angular momentum expressed in its
infinite-dimensional form, related to the $M_{\rho\mu}$ generators.\\
\textbf{Observation:} no spin operator is found, as expected, because the scalar
field belongs to the $(0,0)$ representation, and so the only possible angular momentum
of the system is orbital.

And so the general Lorentz transformation on a scalar field is
$$\phi'(x'^\mu) = \exp\left[-\frac{i}{2}\varepsilon^{\mu\rho}L_{\rho\mu}\right]
\phi(x^\mu)$$

\textbf{Observation:} Any power of a scalar field is a Lorentz scalar (and so all
functions of $\phi$ expressed as power series), the derivative of $\phi$ is
not scalar, while the double derivative is scalar.\\

\textbf{Example:} Lorentz transformation of a vector field.\\
Consider the gradient vector field
$$\pd_\mu \phi,$$
where $\phi$ is a scalar field: what does the operator $\Delta$ do to this kind of
field?
$$\Delta\pd_\mu\phi = \delta\pd_\mu\phi+\delta x^\nu\pd_\nu\pd_\mu\phi$$
since $\Delta \phi =0$, this is also
$$\Delta\pd_\mu\phi = \left[\Delta,\pd_\mu\right]\phi = \left[\delta,\pd_\mu\right]
\phi+\left[\delta x^\nu\pd_\nu,\pd_\mu\right]$$
since $\delta$ is a linear operator, it commutes with the 4-derivative; expanding $
\delta x^\nu$ yields
$$\Delta\pd_\mu\phi = \left[\varepsilon^{\nu\rho}x_\rho\pd_\nu,\pd_\mu\right]=
\varepsilon^{\nu\rho}x_\rho\pd_\nu(\pd_\mu\phi) - \varepsilon^{\nu\rho}\pd_\mu(
x_\rho\pd_\nu\phi)$$
expanding with the product rule
$$\Delta\pd_\mu\phi = \varepsilon^{\nu\rho}x_\rho\pd_\nu\pd_\mu\phi - 
\varepsilon^{\nu\rho}x_\rho\pd_\mu\pd_\nu\phi - \varepsilon^{\nu\rho}(\pd_\mu x_\rho)
\pd_\nu\phi$$
the first two terms are zero, because $\varepsilon^{\nu\rho}\left(x_\rho\pd_\nu\pd_\mu
- x_\rho\pd_\mu\pd_\nu\right)=0$, and $\pd_\mu x_\rho = g_{\mu\rho}$, so
$$\Delta\pd_\mu\phi = -\varepsilon^{\nu\rho}g_{\mu\rho}\pd_\nu\phi =
\varepsilon^{\rho\nu}g_{\mu\rho}\pd_\nu\phi.$$
So finally
\begin{align*}
\delta(\pd_\mu\phi) + \varepsilon^{\rho\sigma}x_\rho\pd_\sigma\pd_\mu\phi = 
\varepsilon\indices{^\rho^\nu}g_{\rho\mu}\pd_\nu\phi\\
\delta(\pd_\mu\phi) = \varepsilon\indices{^\rho^\sigma}g_{\mu\rho}
g\indices{^\nu_\sigma}(\pd_\nu\phi)-\varepsilon^{\rho\sigma}x_\rho\pd_\sigma
(\pd_\mu\phi)
\end{align*}
antisymmetrizing both addends with respect to the indices of $\varepsilon$ yields
$$\delta(\pd_\mu\phi) = \frac{1}{2}\varepsilon^{\rho\sigma}\left[
\underbrace{\left(g_{\mu\rho}g\indices{^\nu_\sigma} - g_{\mu\sigma}
g\indices{^\nu_\rho}\right)}_\text{Independent of $x^\mu$}
(\pd_\nu\phi) - \underbrace{\left(x_\rho\pd_\sigma - x_\sigma\pd_\rho\right)}_\text{
$x^\mu$-dependent}(\pd_\mu\phi)\right].$$
This form is convenient, because it is clear that the transformation of the field has
both a $x^\mu$-dependent generator, which is proportional to the generalized angular
momentum $L_{\rho\sigma}$, and a generator that does not depend from $x^\mu$,
which corresponds to a spin operator $\left(S_{\rho\sigma}\right)\indices{_\mu^\nu}$:
this spin operator is none other than the generator of the finite-dimensional Lorentz
transformations $\left(M_{\rho\sigma}\right)\indices{_\mu^\nu}$ and the orbital
angular momentum operator that of the infinite-dimensional transformations.

$$\delta(\pd_\mu\phi) = -\frac{i}{2}\varepsilon^{\rho\sigma}\left(S_{\rho\sigma}
\right)\indices{_\mu^\nu}(\pd_\nu\phi) -
\frac{i}{2}\varepsilon^{\rho\sigma}L_{\rho\sigma}(\pd_\mu\phi)$$

This example underlines the fact that vector fields can also present internal degrees
of freedom, identified as spin, as they are $(n,m)\neq(0,0)$ irreducible
representations of the Poincaré group.

\lesson{07/10}
\section{Constructing a Lagrangian}
The guiding principle of this process is the Occam's Razor: nature will realize
the simplest Lorentz-invariant terms in the Lagrangian.

The fundamental quantity to build a Lagrangian theory is \textbf{the action}, defined
before as
$$S = \int\dd^4x\mathcal{L}(\phi,\pd_\mu\phi):$$
in classical mechanics the physical trajectory is the one on which the action (as a 
functional) is minimal; in QM, even if no trajectory exists, this process can be still
used in path integral formulation.

As before stated, the Lagrangian density must be \textit{local}, to guarantee
causality, \textit{real}, to preserve probability, \textit{Poincaré 
invariant}, to ensure that the theory is relativistic and must contain \textit{at
most a second derivative}, again to guarantee causality.

Other invariances are admitted (and usually present) and are called \textit{internal
symmetries}: for example the invariance of a spinor under a global phase shift
guarantees \textbf{charge conservation}.
\COMMENT{
\textbf{Example: The Return of Klein-Gordon}\footnote{probably useless,
maybe take it out}\\
The K-G equation can be interpreted as an eigenvalue equation for a field:
$$(\Box+m^2)\phi(x)=0$$
where the d'Alembert operator $\Box = -P_\mu P^\mu$ is a Casimir operator of
eigenvalue $m^2$.
}
\section{Least action}
The action in natural units is dimensionless, so the Lagrangian density must have
units [m$^{-4}$]=[kg$^4$].

Consider an internal symmetry (i.e. a transformation that does not change the
spacetime coordinates)\footnote{this is equivalent to the small variations of the
trajectory in Lagrangian Mechanics.}:
$$\phi(x)\rightarrow \phi'(x')=\phi(x)+\delta\phi(x),$$
the action
$$S=\int_{}^{}\dd^4x\mathcal{L}$$
is minimal if
$$\delta S =0 = \int_{}^{}\dd^4x\delta\mathcal{L} = 
\int_{}^{}\dd^4x \left[\pdv[]{\mathcal{L}}{\phi}\delta\phi + \pdv{\mathcal{L}}
{(\pd_\mu\phi)}\pd_\mu(\delta\phi)\right],
$$
where the relation $\delta(\pd_\mu\phi) = \pd_\mu(\delta\phi)$ was used;
consider
\begin{gather*}
\pd_\mu\left(\pdv[]{\mathcal{L}}{(\pd_\mu\phi)}\delta\phi\right) = 
\pdv[]{\mathcal{L}}{(\pd_\mu\phi)}\pd_\mu(\delta\phi)+\pd_\mu\left(
\pdv[]{\mathcal{L}}{(\pd_\mu\phi)}\right)\delta\phi \rightarrow\\
\rightarrow\pdv[]{\mathcal{L}}{(\pd_\mu\phi)}\pd_\mu(\delta\phi) = \pd_\mu\left(
\pdv[]{\mathcal{L}}{(\pd_\mu\phi)}\delta\phi
\right) - \pd_\mu\left(\pdv[]{\mathcal{L}}{(\pd_\mu\phi)}\right)\delta\phi:
\end{gather*}
substituting it into the $\delta S=0$ condition yields
$$0=\int_{}^{}\dd^4x\left[\pdv[]{\mathcal{L}}{\phi}-\pd_\mu\pdv[]{\mathcal{L}}
{(\pd_\mu\phi)}\right]\delta\phi + \int_{}^{}\dd^4x\pd_\mu\left(\pdv[]{\mathcal{L}}
{(\pd_\mu\phi)}\delta \phi\right).$$
If the variation of the fields $\delta\phi$ goes to zero at the boundary of the
surface enclosing the dominion of integration (i.e. infinity)\footnote{
this is the case for all physical fields.}, the second integral
is zero by the divergence theorem and the action is stationary if and only if
$$\boxed{\pdv[]{\mathcal{L}}{\phi}-\pd_\mu\pdv[]{L}{(\pd_\mu\phi)}=0}$$
which is called the \textbf{Euler-Lagrange equation} for fields.

\textbf{Note}: the E-L equation is invariant under transformations in the form
$$\mathcal{L}'=\mathcal{L}+\pd_\mu\Lambda^\mu$$
where $\Lambda^\mu$ is any 4-vector function: these are called \textit{canonical
transformations}. A trivial example of these transformations are simple translations
of $\mathcal{L}\rightarrow\mathcal{L}+C$.

\section{Noether's Theorem}
This is one of the pivotal theorems of QFT and all of theoretical physics and it
states

\noindent\fbox{%
    \parbox{\textwidth}{%
    \textit{To each symmetry of a system corresponds a 4-current satisfying a
    continuity equation and a conserved quantity, called Noether charge.}
    }%
}

Consider a transformation that acts both on the coordinates and a field
$$\begin{dcases*}
  x'^\mu = x^\mu + \delta x^\mu\\
  \phi'(x')=\phi(x)+\Delta\phi(x)
  \end{dcases*}$$
which does not change the action $S$\footnote{This is the definition of symmetry}.
Staring with $\phi$ as a solution to the E-L equation, the action is invariant under
the transformation if and only if
$$\Delta S =0=\int_{}^{}\dd^4x\Delta\mathcal{L}+\int_{}^{}\delta(\dd^4x)\mathcal{L}$$
where
$$\Delta\mathcal{L} = \underbrace{\pdv[]{\mathcal{L}}{\phi}\delta\phi+
\pdv[]{\mathcal{L}}{\pd_\mu\phi}\pd_\mu\delta\phi}_{
= \delta\mathcal{L}} + \pd_\mu\mathcal{L}\delta x^\mu= \underbrace{\left(\pdv[]{
\mathcal{L}}{\phi} - \pd_\mu\pdv{\mathcal{L}}{\pd_\mu\phi}\right)}_{=0\text{ E-L}}
\delta\phi+\pd_\mu\left(\pdv[]{\mathcal{L}}{\pd_\mu\phi}\delta\phi\right)+
(\pd_\mu\mathcal{L})\delta x^\mu.$$
In this case, since the transformation acts on the coordinates, the variation of
the integration measure must be considered:
$$\delta(\dd^4x)= \dd^4x'-\dd^4x$$
using the Jacobian of the coordinate transformation
$$\dd^4x' = |\det(J(x))|\dd^4x = \left|\det\pdv{x'^\mu}{x^\nu}\right|\dd^4x$$
$$
\delta(\dd^4x)=\dd^4x'-\dd^4x\approx \pd_\mu\delta x^\mu\dd^4x,
$$
since
$$\pdv{x'^\mu}{x^\mu} = \pdv{(x^\mu+\delta x^\mu)}{x^\nu} = g\indices{^\mu_\nu}+
\pd_\nu\delta x^\mu,$$
the property $\det A=e^{\Tr(\ln A)}$ yields
$$\det J(x) \underset{\delta x^\mu\to0}\approx\exp\left[\Tr\pd_\nu\delta x^\mu\right]
\sim 1+\pd_\mu\delta x^\mu.$$
From this follows that the variation of the measure is
$$\delta(\dd^4x) \approx \pd_\mu\delta x^\mu\dd^4x.$$

So the action is stationary if and only if
\begin{gather*}
0 = \Delta S = \int_{}^{}\dd^4x \left[\pd_\mu\left(\pdv[]{\mathcal{L}}{\pd_\mu\phi}
\delta\phi\right)+\underbrace{\pd_\mu\mathcal{L}\delta x^\mu + \mathcal{L}\pd_\mu
\delta x^\mu}_{=\pd_\mu(\mathcal{L}\delta x^\mu)}\right] = \\
= \int_{}^{}\dd^4x \pd_\mu\left(\pdv[]{\mathcal{L}}{\pd_\mu\phi}\delta\phi+
\mathcal{L}\delta x^\mu\right) = 0
\end{gather*}
this happens if and only if
$$\pd_\mu\left(\pdv[]{\mathcal{L}}{\pd_\mu\phi}\delta\phi+
\mathcal{L}\delta x^\mu\right)= 0. $$
The function in parenthesis is \textit{the Noether current}
$$J^\mu =\pdv[]{\mathcal{L}}{\pd_\mu\phi}\delta\phi+\mathcal{L}\delta x^\mu $$
and the found relation is its continuity equation, as
$$\pd_\mu J^\mu = \pd_0 J^0 + \pd_i J^i = \pdv[]{\rho}{t}+\divergence\mathbf{j}=0,$$
where $J^0:=\rho$, the density associated with the 3-current
$\mathbf{j}$.

Integrating the continuity equation yields
$$0=\int_{}^{}\dd^3x \pd_\mu J^\mu = \int_{}^{}\dd^3x \pd_0J^0+
\underbrace{\int_{}^{}\dd^3x\pd_i J^i}_{=0\text{ by div thm}}=
\int_{}^{}d^3x\pdv[]{J^0}{t} = \dv[]{}{t}\int_{}^{}\dd^3xJ^0:=\dv[]{Q}{t}=0$$
so the Noether charge which is conserved is
$$Q = \int_{}^{}d^3xJ^0(x^\mu).$$
It is useful to rewrite the Noether current in terms of $\Delta\phi$:
recalling that
$$\delta\phi = \Delta\phi - \pd_\nu\phi\delta x^\nu$$
the current becomes
$$J^\mu = \underbrace{\left(\mathcal{L}g\indices{^\mu_\nu}-\pdv[]{\mathcal{L}}{
\pd_\mu\phi}\pd_\nu\phi\right)}_{:=-T\indices{^\mu_\nu}}\delta x^\nu +
\pdv[]{\mathcal{L}}{\pd_\mu\phi}\Delta\phi$$
where $T\indices{^\mu_\nu}$ is called the \textit{stress energy tensor}; the sign
is a convention that ensures $T\indices{^0_0}\geq0$.

\end{document}
